\vfill\pagebreak
\section{Repeating: \texttt{loop}, \texttt{break} and \texttt{continue}}
\label{sec:control:loop}

\token{loop} runs a block over and over. When control reaches the closing brace,
it goes back to the opening one, and the block runs again. The \token{loop}
itself never stops, so the block must contain the instruction that ends it:
\token{break}, which leaves the loop at once. The program then continues with
the instruction that follows the loop.

The next listing adds up the numbers typed by the user, as many as there are. It
cannot know in advance how many that is, so it loops, and it breaks out when the
user types \token{0}.

\begin{lstlisting}[style=coloredverbatim, caption=Repeating until the user says stop, label=lst:control:loop]
use std::io;

fn main() {
  let mut total = 0;
  loop {
    let n = read!i32(ask-> "A number (0 to stop): ");
    if n == 0 {
      break;
    }
    total += n;
  }
  println("Total: ", total);
}
\end{lstlisting}
\vspace{-10pt}%
\begin{lstlisting}[style=bashVerb, escapechar=@+]
@+\textcolor{teal!80}{alice@dev:\mtilde}@+$ ./main
A number (0 to stop): 4
A number (0 to stop): 10
A number (0 to stop): 3
A number (0 to stop): 0
Total: 17
\end{lstlisting}

Each time round, the block declares a new \token{n}, which lives until the end of
that repetition, called an \textit{iteration}. \token{total} is declared before
the loop, so it survives from one iteration to the next and accumulates the sum.
Declared inside the block, it would start again at \token{0} on each iteration.

\tipbox{A loop whose \token{break} is never reached runs forever. If a program
  seems stuck, \textit{Ctrl+C} in the terminal stops it.}

\subsection{Skipping the rest of an iteration: \texttt{continue}}

\token{continue} is the other way to jump. Instead of leaving the loop, it
abandons the current iteration and goes straight back to the top of the block
for the next one. The next listing refuses negative numbers: for those, line~12
skips the addition on line~14. It also stops by itself once the total reaches
100, with a second \token{break} on line~16. A loop can have any number of
\token{break}s, and the first one reached ends it.
Figure~\ref{fig:control:loop_jumps} draws the three jumps.

\begin{lstlisting}[style=coloredverbatim, caption=Skipping an iteration, label=lst:control:continue]
use std::io;

fn main() {
  let mut total = 0;
  loop {
    let n = read!i32(ask-> "A number (0 to stop): ");
    if n == 0 {
      break;
    }
    if n < 0 {
      println("Negative numbers are ignored.");
      continue;
    }
    total += n;
    if total >= 100 {
      break;
    }
  }
  println("Total: ", total);
}
\end{lstlisting}
\vspace{-10pt}%
\begin{lstlisting}[style=bashVerb, escapechar=@+]
@+\textcolor{teal!80}{alice@dev:\mtilde}@+$ ./main
A number (0 to stop): 4
A number (0 to stop): -10
Negative numbers are ignored.
A number (0 to stop): 3
A number (0 to stop): 0
Total: 7
@+\textcolor{teal!80}{alice@dev:\mtilde}@+$ ./main
A number (0 to stop): 60
A number (0 to stop): 50
Total: 110
\end{lstlisting}

\begin{figure}[h]
  \centering
  \begin{adjustbox}{max width=\linewidth}
    \begin{tikzpicture}[node distance=6mm]

      \node[cfterm] (start) {start of main};
      \node[cfstep, below=of start] (init) {let mut total = 0;};
      \node[cfstep, below=8mm of init] (read) {let n = read!i32(...);};
      \node[cftest, below=of read] (zero) {n == 0};
      \node[cftest, below=of zero] (neg) {n < 0};
      \node[cfstep, below=of neg] (add) {total += n;};
      \node[cftest, below=of add] (full) {total >= 100};
      \node[cfstep, below=8mm of full] (print) {println("Total: ", total);};
      \node[cfterm, below=of print] (end) {end of main};
      \node[cfstep, right=12mm of neg] (ignored) {println("... ignored.");};

      \coordinate (out) at ($(full.west) + (-12mm, 0)$);
      \coordinate (back) at ($(ignored.east) + (8mm, 0)$);

      \draw[cfflow] (start) -- (init);
      \draw[cfflow] (init) -- node[cflab, right] {loop \{} (read);
      \draw[cfflow] (read) -- (zero);
      \draw[cfflow] (zero) -- node[cflab, right] {false} (neg);
      \draw[cfflow] (neg) -- node[cflab, right] {false} (add);
      \draw[cfflow] (add) -- (full);
      \draw[cfflow] (print) -- (end);

      \draw[cfflow] (zero.west) -- node[cflab, above] {true} (zero.west -| out)
        -- node[cfjump, left, pos=0.5] {break} (out) |- (print.west);
      \draw[cfflow] (full.west) -- node[cflab, above] {true} (out) |- (print.west);

      \draw[cfflow] (neg.east) -- node[cflab, above] {true} (ignored.west);
      \draw[cfflow] (ignored.north) -- node[cfjump, left] {continue} (ignored.north |- read.east)
        -- (read.east);
      \draw[cfflow] (full.east) -- node[cflab, above] {false} (full.east -| back)
        -- node[cflab, right, align=left] {end of\\the block} (back |- read.east)
        -- (read.east);
    \end{tikzpicture}
  \end{adjustbox}
  \caption{\label{fig:control:loop_jumps} The \token{loop} of
    Listing~\ref{lst:control:continue}. The end of the block and
    \token{continue} both go back to the top of the block; either
    \token{break} leaves the loop.}
\end{figure}


\token{break} and \token{continue} both make the lines that follow them in the
same block unreachable, since control never gets there. In the next listing,
the goodbye message on line~9 can never be printed: the \token{break} just
above it has already left the loop.

\begin{lstlisting}[style=coloredverbatimError, escapechar=@, caption=An instruction after a \token{break}]
use std::io;

fn main() {
  let mut total = 0;
  loop {
    let n = read!i32(ask-> "A number (0 to stop): ");
    if n == 0 {
      break;
      @\hb{println("Goodbye.")}@;
    }
    total += n;
  }
  println("Total: ", total);
}
\end{lstlisting}

The compiler rejects such lines with ``unreachable statement'', as it rejects a
condition that is always true: code that can never run is a mistake. Here the
message belongs before the \token{break}. For the same reason, both words are
errors outside of a loop, where there is nothing to leave and nothing to repeat.

\subsection{A loop that produces a value}

A \token{loop}, like an \token{if}, can be used as a value. The value comes from
the \token{break} that ends it: \token{break m;} leaves the loop, and \token{m}
becomes the value of the whole \token{loop}.

The next listing uses this to insist until the user gives a valid month. The
loop only ends through the \token{break} on line~7, which carries the month out,
so \token{month} is always between 1 and 12.

\begin{lstlisting}[style=coloredverbatim, caption=A loop whose value comes from \token{break}, label=lst:control:loop_value]
use std::io;

fn main() {
  let month = loop {
    let m = read!i32(ask-> "Month (1 to 12): ");
    if m >= 1 && m <= 12 {
      break m;
    }
    println("There is no month ", m, ".");
  };
  println("Month ", month, " it is.");
}
\end{lstlisting}
\vspace{-10pt}%
\begin{lstlisting}[style=bashVerb, escapechar=@+]
@+\textcolor{teal!80}{alice@dev:\mtilde}@+$ ./main
Month (1 to 12): 0
There is no month 0.
Month (1 to 12): 13
There is no month 13.
Month (1 to 12): 4
Month 4 it is.
\end{lstlisting}

The semicolon after the closing brace on line~10 ends the declaration of
\token{month}, as it would after any other value.

A \token{loop} with several \token{break}s takes its value from whichever one
ends it, so they must all carry a value, of the same type. The next listing
lets the user cancel by typing \token{0}, but its new \token{break} on line~7
carries nothing.

\begin{lstlisting}[style=coloredverbatimError, escapechar=@, caption=Two \token{break}s that do not agree]
use std::io;

fn main() {
  let month = loop {
    let m = read!i32(ask-> "Month (1 to 12, 0 to cancel): ");
    if m == 0 {
      break;
    }
    if m >= 1 && m <= 12 {
      @\hb{break m}@;
    }
    println("There is no month ", m, ".");
  };
  println("Month ", month, " it is.");
}
\end{lstlisting}

The compiler reports ``incompatible types void and i32'' on the second
\token{break}: the first one gave the loop no value, of type \token{void}, and
the second an \token{i32}. The cancellation needs a value of its own, such as
\token{break 0;}, which the rest of the program can then test.
